\begin{textsample}{POS Dim 6 – ac – Score 9.00 – ac/ac\_em\_30.txt}  \label{ex:f6_pos_153}
14.
Eixos \textbf{estruturantes} A organização das Rotas de Aprofundamento se dá conforme orienta o parágrafo 2º do Artigo 12 das DCNEM, que estabelece quatro eixos \textbf{estruturantes} complementares, estes : Investigação Científica, Processos Criativos, \textbf{Mediação} e Intervenção Sociocultural e \textbf{Empreendedorismo}.
Tais eixos visam garantir aos estudantes oportunidades de experimentar diferentes situações de aprendizagem, desenvolvendo um conjunto diversificado de habilidades relevantes para sua formação integral.
Desse modo, no estado do Acre optou -se por incorporar os quatro eixos, ao longo do desenvolvimento do itinerário \textbf{formativo} do aluno, independente da área de aprofundamento escolhida pelo estudante.
De acordo com o Parecer do MEC nº 1.432/2018, são estabelecidos justificativa, objetivo e foco pedagógico para cada um dos eixos, além de habilidades gerais de cada eixo e habilidades específicas dos eixos para cada área de conhecimento.
Assim, segundo os referenciais para a elaboração de itinerários \textbf{formativos}, temos : - Investigação Científica : centra -se na realização de práticas e produções científicas.
Justifica -se pela necessidade social de apropriação de conhecimentos e habilidades que permitam acessar, selecionar, \textbf{processar}, analisar e utilizar dados sobre os mais diferentes assuntos.
Assim, tem como objetivo aprofundar conhecimentos fundantes das ciências ; ampliar as habilidades de pensar e fazer científico ; e compreender e enfrentar situações cotidianas, intervindo na sociedade.
O foco pedagógico é a pesquisa científica. - Processos Criativos : centra -se na idealização e realização de projetos criativos.
Justifica -se pela necessidade de participação em uma sociedade cada vez mais pautada pela criatividade e inovação.
Tendo como objetivo aprofundar os conhecimentos sobre as artes, a cultura, as mídias e as ciências aplicadas ; ampliar habilidades relacionadas ao pensar e fazer criativo ; e criar e produzir soluções inovadoras e criativas para \textbf{problemas} identificados na sociedade.
O foco pedagógico é a realização de projetos criativos. - \textbf{Mediação} e intervenção sociocultural : centra -se na realização de projetos que contribuam com a sociedade e o meio ambiente.
Justifica -se pela necessidade de formação voltada para a atuação como agentes de mudanças e de construção de uma sociedade mais ética, justa, democrática, inclusiva, solidária e sustentável.
Assim, tem -se como objetivo aprofundar conhecimentos sobre questões que afetam a vida dos seres humanos e do planeta ; ampliar as habilidades relacionadas à convivência e atuação sociocultural ; e desenvolver habilidades de mediar conflitos, promover entendimentos e propor soluções para questões e \textbf{problemas} socioculturais e \textbf{ambientais} identificados em suas comunidades.
O foco pedagógico é o engajamento em projetos de mobilização e intervenção sociocultural e \textbf{ambiental}. - \textbf{Empreendedorismo} : centra -se em empreender projetos pessoais ou produtivos articulados ao seu projeto de vida.
Justifica -se pela necessidade de adaptabilidade a uma sociedade cada vez mais volátil que exige um \textbf{perfil} que se adapte aos diferentes contextos e crie novas oportunidades para si e para os demais.
Objetiva -se aprofundar conhecimentos relacionados ao contexto, ao mundo do trabalho e à gestão de iniciativas empreendedoras ; ampliar habilidades relacionadas ao autoconhecimento, \textbf{empreendedorismo} e projeto de vida ; e estruturar iniciativas empreendedoras com propósitos diversos.
O foco pedagógico é a criação de empreendimentos pessoais ou produtivos articulados com seus projetos de vida.

% matched lemmas: ambiental, empreendedorismo, estruturante, formativo, mediação, perfil, problema, processar
\end{textsample}
